\documentclass[11pt]{article}


\setlength{\parindent}{0pt}
\usepackage{xltxtra}
\usepackage{hyperref}
\hypersetup{hidelinks}
\usepackage{url}
\urlstyle{tt}
\usepackage{xcolor}
\definecolor{CVBlue}{RGB}{23,110,191}
\usepackage{calc}
\usepackage{graphicx}
\usepackage{tikz}
\usetikzlibrary{calc}
\usepackage{fontspec}
\usepackage{xeCJK}
\usepackage{enumitem}
\CJKsetecglue{} %% 取消中文与数字之间的间隙


%% 主文档字体设置
\setmainfont[
    Path = fonts/Main/,
    Extension = .otf,
    BoldFont = texgyretermes-bold.otf, % 加粗字体
]{texgyretermes-regular.otf} % 正文字体

% 中文字体设置
\setCJKmainfont[
    Path = fonts/hansans/,
    Extension = .ttf,
    BoldFont = NotoSansSC-Bold.ttf, % 加粗字体
]{NotoSansSC-Regular.otf} % 正文字体


\usepackage{relsize}
\usepackage{xspace}

% 使用 fontawesome（部分图标）
\usepackage{fontawesome} 

% A4纸，上下左右边距
\usepackage[
    a4paper,
    left=1.2cm,
    right=1.2cm,
    top=1.5cm,
    bottom=1cm,
    nohead
]{geometry}

\renewcommand{\baselinestretch}{1.5} % 行间距设为1.5

\usepackage{titlesec}
\usepackage{enumitem}
\setlist{noitemsep} % 取消列表项间的额外间距
%\setlist{nosep} % 取消所有垂直间距
\setlist[itemize]{topsep=0.25em, leftmargin=*}
\setlist[enumerate]{topsep=0.25em, leftmargin=*}

% --- 用于控制【不同项目之间】的垂直距离 ---
\newlength{\interProjectSpacing}
\setlength{\interProjectSpacing}{0.9em} % <--- 在此调整项目之间的距离
\newcommand{\projectsep}{\vspace{\interProjectSpacing}}

% --- 用于控制【项目标题】与下方【项目描述】的距离 ---
\newlength{\intraProjectTitleSep}
\setlength{\intraProjectTitleSep}{0.4em} % <--- 在此调整标题和描述的距离
\newcommand{\titlebreak}{\\[\intraProjectTitleSep]}

% --- 用于控制【项目描述】与下方【要点列表】的距离 ---
\newlength{\intraProjectListTopSep}
\setlength{\intraProjectListTopSep}{0.2em} % <--- 在此调整描述和列表的距离

% =======================================================================


\titleformat{\section}         % 定制 \section 命令 
{\large\bfseries\raggedright} % 将 section 标题设置为大号、粗体且左对齐
{}{0em}                      % 可用于为所有 section 添加前缀（如“章节...”）
{}                           % 可用于在标题前插入代码
[{\color{CVBlue}\titlerule}]  % 在标题后插入一条横线
\titlespacing*{\section}{0cm}{*1.6}{*1.2}



\begin{document}
\pagenumbering{gobble}

%%%% 利用tikz来定位照片
\begin{tikzpicture}[remember picture, overlay] 
    % \node[anchor = north east] at ($(current page.north east)+(-2cm,-1.2cm)$) {\includegraphics[height=3cm]{avatar.jpg}};
  \end{tikzpicture}%
  %%%% 利用tikz来定位学校Logo，这里只在第一页显示
  \begin{tikzpicture}[remember picture, overlay] 
    \node[anchor = north west] at ($(current page.north west)+(0.5cm,+1.0cm)$) {\includegraphics[height=6cm]{zju.png}};
  \end{tikzpicture}%
\centerline{\LARGE\bfseries{陈一天}} 

\centerline{\normalsize{\faPhone\ 189-8011-1110 \quad \faEnvelopeO\ \href{mailto:chenyitian01@zju.edu.cn}{chenyitian@zju.edu.cn}}} 

\centerline{\normalsize{\faGithubSquare\ \href{https://github.com/yitian-chen}{https://github.com/yitian-chen}}} 
    
\section{\makebox[\widthof{\faGraduationCap}][c]{\color{CVBlue}\faGraduationCap}\ 教育背景}    
\textbf{浙江大学} \hfill 电子科学与技术 \hfill 2024.9 -- 2028.6\\[0.5em]
% \begin{itemize}[nosep]
%     \item 相关课程：
% \end{itemize}

\section{\makebox[\widthof{\faUsers}][c]{\color{CVBlue}\faUsers}\ 项目经历}

% --- 第一个项目 ---
% 将标题行末尾的 \\ 替换为 \titlebreak 命令
\textbf{LeaseGo | 公寓租赁平台} \hfill 2026.02 -- 2026.04 \titlebreak
{\small\color{CVBlue}仓库链接 \faGithubSquare\ \href{https://github.com/yitian-chen/LeaseGo}{https://github.com/yitian-chen/LeaseGo}}\\*
\textbf{技术栈}: Java 17 + Spring Boot 3.2 + Spring Cloud Alibaba + MyBatis-Plus + MySQL 8.0, Redis Stack + RabbitMQ + Redisson + WebSocket + MinIO + LangChain4j + Docker
% 在 itemize 的选项中，使用 topsep=\intraProjectListTopSep 来控制上边距
\begin{itemize}[nosep, topsep=\intraProjectListTopSep]
    \item 独立设计并实现基于 \textbf{Spring Cloud Alibaba} 的微服务架构，将app端拆分为 \textbf{Gateway} 网关、\textbf{web-app} 用户端、\textbf{chat-service} 即时通讯、\textbf{agent-service} AI 检索四个独立服务，采用 \textbf{Nacos} 实现服务发现与配置中心管理。

    \item 基于 \textbf{LangChain4j} + \textbf{MiniMax LLM} + 阿里 \textbf{DashScope Embedding} 构建 \textbf{RAG} 检索增强生成管道，将结构化房源数据嵌入 \textbf{Redis Stack} 向量库，实现自然语言驱动的智能房源推荐。

    \item 集成 \textbf{RabbitMQ} 异步消息队列，实现房间变更自动重索引、聊天消息持久化、浏览历史记录、租约到期通知等四项异步任务，采用手动确认 + 重试机制保障消息可靠投递。

    \item 基于 \textbf{WebSocket} + \textbf{Redis Pub/Sub} 实现跨实例实时通讯，支持多实例部署下的聊天消息分发与已读未读状态追踪，打通房源详情页与房东即时沟通链路。

    \item 使用 \textbf{Redisson} 分布式锁按房间粒度防护租约签约超卖、进行聊天会话去重。

    \item 封装 \textbf{ThreadLocal} 上下文持有者 + 自定义拦截器实现 \textbf{JWT} 无感认证，建立统一全局异常拦截器与标准 \textbf{Result} 返回格式。
\end{itemize}

% 使用 \projectsep 命令来分隔两个项目
\projectsep

% --- 第二个项目 ---
\textbf{Synchro | AI智能匹配交友平台} \hfill 2026.04 -- 2026.05 \titlebreak
{\small\color{CVBlue}仓库链接 \faGithubSquare\ \href{https://github.com/yitian-chen/Synchro}{https://github.com/yitian-chen/Synchro}}\\*
\textbf{技术栈}Java 17 + Spring Boot 3.3 + LangChain4j + DashScope Embedding + MySQL 8 + JPA + Flyway + Redis Stack + STOMP WebSocket + MinIO + Docker
\begin{itemize}[nosep, topsep=\intraProjectListTopSep]                                
    \item 基于 \textbf{LangChain4j} 与大语言模型实现 8 轮 \textbf{AI 性格面试}，从自然对话中提取 22 维人格特征，替代传统问卷，提升用户体验与数据密度。
    
    \item 设计\textbf{多因子加权匹配算法}，融合特征余弦相似度、1536 维语义嵌入相似度、互补性格匹配及双向理想画像匹配，实现自动化周度配对。
    
    \item 使用 \textbf{Redis Stack} 作为向量存储与缓存层，支持高性能语义检索与实时消息缓存。
    
    \item 基于 \textbf{STOMP over WebSocket} 实现实时聊天。
    
    % \item 采用 Resilience4j 熔断重试、Flyway 数据库版本管理、MinIO 对象存储等工程化实践                                 
\end{itemize}


\section{\makebox[\widthof{\faCogs}][c]{\color{CVBlue}\faCogs}\ 专业技能}
\begin{itemize}[nosep]
    \item \textbf{Java：} 具有扎实的Java基础，具备良好的面向对象的编程思想。
    \item \textbf{MySQL：} 熟悉SQL语言，熟悉事务及其原理、存储引擎、索引、锁机制等。
    \item \textbf{Redis：} 熟悉Redis的基本使用，了解其数据类型、缓存过期策略、淘汰策略、持久化机制。
    \item \textbf{框架：} 熟练运用Springboot开发框架，理解Spring IOC、AOP、Bean生命周期
    \item \textbf{其他：} 熟练使用 Claude Code 等 Vibe Coding 工具实现高效开发与调试；熟悉 IDEA 等开发工具、Git 版本管理、Docker 部署
\end{itemize}


% \section{\makebox[\widthof{\faGraduationCap}][c]{\color{CVBlue}\faList}\ 获奖情况}
% \begin{itemize}
%     \item CC98年度用户 \hfill 2023.12
    
% \end{itemize}
    
% \section{\makebox[\widthof{\faInfo}][c]{\color{CVBlue}\faInfo}\ 其他}
% \begin{itemize}[parsep=0.5ex]
%     \item \textbf{技术博客：} \href{https://maksymilan.github.io/}{https://maksymilan.github.io/}
%     \item \textbf{GitHub：} \href{https://github.com/maksymilan}{https://github.com/maksymilan} 
%     \item \textbf{英语水平：} CET-4, CET-6
% \end{itemize}
\end{document}